A rectifiable path to infinity is understood to be locally rectifiable; its total length is necessarily infinite. For an entire function $f$, let $\ell(r)$ be the length of its asymptotic path inside $|z|<r$. We construct a nonconstant entire function of order zero such that
\[
 \log M(r,f)=O((\log r)^3)
 \quad\text{and}\quad
 \limsup_{r\to\infty}\frac{\ell(r)}r=\infty
\]
for every path on which $f(z)\to\infty$. Thus finite order does not guarantee a path with $\ell(r)=O(r)$.

This is a reconstruction of the spiral-obstacle method of Gol'dberg--Eremenko, \emph{On asymptotic curves of entire functions of finite order}, Section 1, \url{https://www.math.purdue.edu/~eremenko/dvi/as-curves.pdf}. Their sharper result permits any prescribed divergent factor in place of the extra $\log r$ above. We prove the displayed, sufficient order-zero counterexample and do not claim that sharper growth estimate.
